\subsection{Spectral Stability and the Unit of Mass}
  \label{sec:spectral-stability}

  Within the admissibility framework developed in Chapter~\ref{sec:spectral-admissibility},
  rest mass appears as a spectral invariant associated with
  the stability of projected configurations.

  In the effective description, particle states correspond to
  stable modes supported by the projection fiber~$\Pi$.
  Their mass is determined by the eigenvalues of the scalar
  Laplacian~$\Delta^{(0)}_G$ acting on this space, subject to
  topological constraints~$\mathcal{T}$:
  \begin{equation}
    m^2 c^2 \;=\; \lambda_{\mathcal{T}}
    \;\equiv\;
    \mathrm{Eig}\!\left(\Delta^{(0)}_G\right)
    \Big|_{\mathcal{T}} .
    \label{eq:mass_spectral_def}
  \end{equation}

  The quantity $\lambda_{\mathcal{T}}$ represents the minimal
  spectral cost required to maintain the corresponding
  configuration within the admissible projective sector.
  Mass therefore measures the resistance of a configuration
  to projective ordering advancement.

  The reference mass scale~$m_{(1)}$ corresponds to the lowest
  non-trivial admissible eigenmode~$\lambda_1$ of the fiber
  geometry $\Pi \cong S^3$.
  All other sector mass scales can then be expressed relative to
  this fundamental scale,
  \begin{equation}
    m \;=\; m_{(1)}
    \sqrt{\frac{\lambda_{\mathcal{T}}}{\lambda_1}} .
  \end{equation}
